\clearpage\begingroup\SourceLeaf{72b}\HistoricalHead{67}{R}{BERECHNUNG MAGNETISCHER KRÄFTE.}\HistoricalFoot{}{9\textsuperscript{*}}
\begin{center}[9.]\end{center}
\noindent{\small [S. 217]}
\begin{center}\textbf{Berechnung gedämpfter Bewegungen der Magnetnadel.}\end{center}
\noindent\hspace*{\parindent}\begin{tabular}{@{}l@{}}
\(T\) Schwingungsdauer (vergrösserte)\\
\(a:b\) Abnahme-Verhältniss der Bogen\\
\(t,t+T\) Beobachtungszeiten\\
\(A,A+\Delta\) Stand\\
\(k\) Modulus der brigg. Log.
\end{tabular}\par
\indent Dann gilt der Stand \(A+\dfrac{a}{a+b}\Delta\) für die Zeit
\[
t+\left(\frac{a}{a+b}-\frac{2\log\left(\dfrac ab\right)}{k\left(\pi^2+\left[\dfrac1{k^2}\right]\left(\log\dfrac ab\right)^2\right)}\right)
\]
\noindent oder, wenn man
\[
\frac{\log\dfrac ab}{k\pi}=\operatorname{tg}\varphi
\]
\clearpage\SourceLeaf{73a}\HistoricalHead{68}{L}{PHYSIK. NACHLASS.}
\noindent setzt, wo dann \(T\cos\varphi\) die ungedämpfte Schwingungsdauer wird,
\[
t+\left(\frac{a}{a+b}-\frac{\sin2\varphi}{\pi}\right)T.
\]
\noindent[Hier folgen Zahlenbeispiele.]
\bigskip\begin{center}\textbf{BEMERKUNG.}\end{center}
{\small\noindent\hspace*{\parindent}Der Zweck der obigen Betrachtung geht aus einer Vergleichung mit der Abhandlung \textit{Anleitung zur
Bestimmung der Schwingungsdauer einer Magnetnadel} hervor, man vergleiche namentlich Werke V, S. 393.
\par\raggedleft\textsc{Schaefer}.\par}
\clearpage\endgroup
